\section*{Chapter \ref{ch:g10:elasticities} --- Elasticities}
% figdata: figdata/grade-10/07-elasticities/out/ (health, health_fit, health_marked, food_fit); ledger g10:gasoline-elasticity,
% g10:cigarette-elasticity; Brenland numbers invented.

\begin{solution}{exo:g10:elasticities:1}
$\bar Q = 550$, $\bar p = 4.5$: $\varepsilon_d = (-100/550)/(1/4.5) = -0.1818/0.2222 \approx -0.82$. Its size is below 1: demand is inelastic in this
range.
\end{solution}

\begin{solution}{exo:g10:elasticities:2}
$\bar Q = 230$, $\bar p = 9$: $\varepsilon_d = (60/230)/(-2/9) = 0.2609/(-0.2222) \approx -1.17$, elastic. Revenue rises, from $10 \times 200 = 2000$ to $8
\times 260 = 2080$ dalens a day: with elastic demand, a lower price brings more revenue.
\end{solution}

\begin{solution}{exo:g10:elasticities:3}
Income: $65228/6945 \approx 9.4$ times. Health spending: $10550/211 = 50$ times. The fitted slope of 1.12 would give about $9.4^{1.12} \approx 12$ times: the
United States lies well above the line (and India somewhat below it).
\end{solution}

\begin{solution}{exo:g10:elasticities:4}
(a) Positive (\omterm{def:g10:demand:substitutes}{substitutes}). (b) Negative (\omterm{def:g10:demand:substitutes}{complements}). (c) Negative (\omterm{def:g10:demand:substitutes}{complements}). (d) Positive (\omterm{def:g10:demand:substitutes}{substitutes}).
\end{solution}

\begin{solution}{exo:g10:elasticities:5}
Inelastic; unit-elastic; elastic; \omterm{def:g10:elasticities:perfectly-elastic}{perfectly inelastic}; \omterm{def:g10:elasticities:perfectly-elastic}{perfectly elastic}.
\end{solution}

\begin{solution}{exo:g10:elasticities:6}
At $p = 10$: $Q = 150$, $\varepsilon_d = -5 \times 10/150 \approx -0.33$. At $p = 20$: $Q = 100$, $\varepsilon_d = -100/100 = -1$. At $p = 30$: $Q = 50$,
$\varepsilon_d = -150/50 = -3$. Unit-elastic at $p = 200/(2 \times 5) = 20$, the midpoint of the line.
\end{solution}

\begin{solution}{exo:g10:elasticities:7}
At $p = 10$: $Q = 20$, $\varepsilon_s = 4 \times 10/20 = 2$. At $p = 20$: $Q = 60$, $\varepsilon_s = 80/60 \approx 1.33$.
\end{solution}

\begin{solution}{exo:g10:elasticities:8}
By \cref{prop:g10:elasticities:revenue}, $p = 90/(2 \times 3) = 15$, $Q = 90 - 45 = 45$ and $R = 15 \times 45 = 675$. There $\varepsilon_d = -3 \times 15/45 =
-1$.
\end{solution}

\begin{solution}{exo:g10:elasticities:9}
The poorest fifth spends $3940/13870 \approx 28.4\,\%$ of its income on food, and the richest fifth $5340/41790 \approx
12.8\,\%$. Arc income elasticity: $(1400/4640)/(27920/27830) \approx 0.302/1.003 \approx 0.30$. If
the richest fifth had spent 28.4\,\% too, food spending would be proportional to income, and the \omterm{def:g10:elasticities:arc}{arc elasticity} would be exactly 1: the two percentage
changes, each measured against its own mean, would be equal.
\end{solution}

\begin{solution}{exo:g10:elasticities:10}
The quantity falls by about $0.4 \times 10 = 4\,\%$, to about $4800$ tickets. Revenue goes from $10 \times 5000 = \text{\dal{50000}}$ to $11 \times 4800 = \text{\dal{52800}}$,
a rise of \dal{2800} (5.6\,\%): with inelastic demand, the higher price brings more revenue.
\end{solution}

\begin{solution}{exo:g10:elasticities:11}
(a) About $0.09 \times 10 \approx 0.9\,\%$ within a year and $0.31 \times 10 \approx 3.1\,\%$ in the long run. (b) Adults: between $0.3 \times 20 = 6\,\%$ and $0.6
\times 20 = 12\,\%$. Young people: between 12\,\% and 14\,\%.
\end{solution}

\begin{solution}{exo:g10:elasticities:12}
A rise in the price, \omterm{def:g10:economic-models:other-things-equal}{other things equal}, after which the quantity fell by a larger percentage than the price rose, so that buyers' spending on the good fell.
Check that the price change came from the supply side (a tax, a cost) and not from a shift of demand; that incomes, the prices of \omterm{def:g10:demand:substitutes}{substitutes} and the number
of buyers did not change at the same time; that buyers had not stocked up before the rise; and over what time span the response was measured, since
elasticities are larger in the long run.
\end{solution}

\begin{solution}{exo:g10:elasticities:13}
$\varepsilon_d = -bp/(a - bp)$; dividing top and bottom by $b$ gives $-p/(a/b - p) = -p/(P - p)$. It equals $-1$ when $p = P - p$, that is $p = P/2$. At $p =
3P/4$: $\varepsilon_d = -(3P/4)/(P/4) = -3$.
\end{solution}

\begin{solution}{exo:g10:elasticities:14}
Spending is $p\,Q = 1200$ at every price. From $p = 2$ to $p = 3$, $Q$ goes from 600 to 400: $(-200/500)/(1/2.5) = -0.4/0.4 = -1$. A curve with the same elasticity
everywhere is one on which every 1\,\% rise in the price cuts the quantity by the same percentage, wherever one starts. A straight line cannot do this: by
\cref{prop:g10:elasticities:midpoint} its elasticity rises from 0 at the bottom to very large sizes at the top. A constant-elasticity curve bends: it is steep
at small quantities, flat at large ones, and never meets either axis.
\end{solution}

\begin{solution}{exo:g10:elasticities:15}
In the short run, drivers own the cars they own, live where they live and work where they work: they can only skip a few trips or share a car. Over years they
replace their car by a smaller one, move or change jobs, and car makers build more frugal models: the same price rise cuts fuel use far more (0.31 against 0.09
in the table). For supply, \cref{ex:g10:supply:two-curves} gives $0.4$ in the short run and $1.6$ in the long run: in the short run firms can only use their
capacity harder, while in the long run they build new capacity and new firms enter.
\end{solution}

\begin{solution}{pb:g10:elasticities:1}
Tolls in dalens; crossings a day; revenue in dalens a day.

\textbf{1.} $b = (9000 - 8000)/(5 - 4) = 1000$ crossings per dalen, and $a = 9000 + 4 \times 1000 = 13000$: $Q = 13000 - 1000\,p$.

\textbf{2.} Nobody crosses at $p = 13$; with no toll, \num{13000} crossings.

\textbf{3.} $\varepsilon_d = -1000 \times 5/8000 = -0.625$.

\textbf{4.} $(-1000/8500)/(1/4.5) = -0.1176/0.2222 \approx -0.53$.

\textbf{5.} $4 \times 9000 = \text{\dal{36000}}$ and $5 \times 8000 = \text{\dal{40000}}$. Demand is inelastic in this range, so the higher toll raised revenue, as
\cref{prop:g10:elasticities:revenue} says.

\textbf{6.} $R(p) = p\,(13000 - 1000p) = 13000p - 1000p^2$.

\textbf{7.} $p = 13000/2000 = 6.5$ dalens, with $Q = 13000 - 6500 = 6500$ crossings.

\textbf{8.} $R = 6.5 \times 6500 = \text{\dal{42250}}$; $\varepsilon_d = -1000 \times 6.5/6500 = -1$.

\textbf{9.} $6 \times 7000 = \text{\dal{42000}}$ and $7 \times 6000 = \text{\dal{42000}}$: both below \dal{42250}, symmetric around 6.5.

\textbf{10.} $8 \times 5000 = \text{\dal{40000}}$. Above 6.5 dalens demand is elastic: each rise loses a larger share of crossings than it adds to the toll.

\textbf{11.} $12000 - 6000 = 6000$ crossings instead of \num{8000}: \num{2000} lost.

\textbf{12.} $\varepsilon_d = -1200 \times 5/6000 = -1$.

\textbf{13.} $p = 12000/2400 = 5$ dalens, $Q = 6000$, $R = \text{\dal{30000}}$.

\textbf{14.} $Q = 12000 - 1200 \times 6.5 = 4200$ and $R = 6.5 \times 4200 = \text{\dal{27300}}$, \dal{2700} less than at 5 dalens.

\textbf{15.} The ferry is a close substitute: when the toll rises, travellers now have somewhere else to go, so they leave the bridge in larger numbers.

\textbf{16.} $10800 - 1200p + 400 \times 3 = 12000 - 1200p$.

\textbf{17.} $\varepsilon_{XY} = 400 \times 3/6000 = 0.2$. Positive: bridge and ferry are \omterm{def:g10:demand:substitutes}{substitutes}.

\textbf{18.} $Q = 10800 - 6000 + 1600 = 6400$ crossings; revenue $5 \times 6400 = \text{\dal{32000}}$.

\textbf{19.} $\varepsilon_I \approx 3/10 = 0.3$. Positive and below 1: a \omterm{def:g10:demand:normal}{normal good} and a necessity for these travellers.

\textbf{20.} Before the ferry, the revenue-maximising toll is 6.5 dalens and brings \dal{42250} a day; after the ferry it is 5 dalens and brings \dal{30000} a
day.
\end{solution}
